\section{Experiments}

\subsection{Experimental Setup}

\textbf{Dataset.} We evaluate on FreeGraspData~\cite{freegrasp}. The evaluation set contains three language splits. Each split has six strata: Easy, Medium, and Hard scenes, each with and without ambiguity. We use all 900 instructions.

\textbf{Metrics.} RSR measures whether the selected first-grasp object ID is acceptable under the benchmark grasp-order annotation. SSR measures whether the predicted first-grasp mask reaches IoU $\geq 0.5$ with an acceptable ground-truth mask. We report both metrics because correct reasoning can still fail at segmentation, and an accidental mask overlap should not be treated as a correct target/blocker decision.

\textbf{Baselines.} We compare against the original FreeGrasp GPT-4o results reported by Jiao \emph{et al.}~\cite{freegrasp}. We also include S0, a controlled same-cache baseline that runs the original FreeGrasp protocol with the same MiMo-v2.5 backbone and the same reused Molmo cache as EGHC3\_v3. Finally, S1 is reported only as an oracle upper bound because it uses ground-truth-derived features and benchmark difficulty labels.

\textbf{Implementation.} EGHC3\_v3 uses MiMo-v2.5 as the main VLM, a route guard recorded as \texttt{eghc\_consensus\_v1}, and \texttt{eghc\_mask\_repair\_v1} for segmentation repair. The run uses the preserved Molmo cache recorded in the run manifest (SHA-256 prefix \texttt{705d0883}). The full run manifest records 1889 reasoning calls, 900 router calls, 59 repair calls, and 2.10 calls per sample.

\subsection{Main Results}

Table~\ref{tab:main} shows the headline comparison. EGHC3\_v3 achieves 60.89\% RSR and 55.22\% SSR over 900 instructions. Relative to the reported original FreeGrasp GPT-4o baseline, this corresponds to approximately +13.6 percentage points RSR and +17.0 percentage points SSR. Relative to the controlled S0 original-protocol baseline, EGHC3\_v3 improves RSR by +3.00 percentage points and SSR by +7.67 percentage points.

\begin{table}[t]
\centering
\caption{Main FreeGraspData results. The original FreeGrasp row is the reported GPT-4o baseline from the prior paper; S0 and EGHC3\_v3 are controlled runs under the same reused Molmo cache.}
\label{tab:main}
\small
\begin{tabular}{lcccc}
\toprule
Method & GT features & Difficulty label & RSR & SSR \\
\midrule
FreeGrasp GPT-4o~\cite{freegrasp} & No & No & 47.3 & 38.2 \\
S0 original protocol & No & No & 57.89 & 47.56 \\
\textbf{EGHC3\_v3 (ours)} & \textbf{No} & \textbf{No} & \textbf{60.89} & \textbf{55.22} \\
\bottomrule
\end{tabular}
\end{table}

The larger gain in SSR than RSR indicates that EGHC does not only change the selected first-grasp object; it also improves the consistency between the chosen candidate and the final segmentation mask. This is important for robotic execution, where a correct target/blocker decision still requires an actionable mask.

\subsection{Oracle Headroom}

Table~\ref{tab:oracle} separates the deployable method from the oracle diagnostic result. S1 reaches 67.44\% RSR and 60.11\% SSR, but it uses annotation-derived features and benchmark difficulty labels. We therefore interpret it as an upper-bound-style diagnostic, not as an automatic method. The gap between S1 and EGHC3\_v3 is 6.55 percentage points RSR and 4.89 percentage points SSR, suggesting that more accurate predicted visibility and blocker evidence could still improve the deployable pipeline.

\begin{table}[t]
\centering
\caption{Oracle upper bound. S1 is diagnostic only because it uses ground-truth features and benchmark difficulty labels during inference.}
\label{tab:oracle}
\small
\begin{tabular}{lcccc}
\toprule
Method & GT features & Difficulty label & RSR & SSR \\
\midrule
EGHC3\_v3 (deployable evidence) & No & No & 60.89 & 55.22 \\
S1 oracle diagnostic & Yes & Yes & 67.44 & 60.11 \\
\bottomrule
\end{tabular}
\end{table}

\subsection{Per-Stratum Analysis}

Table~\ref{tab:strata} breaks EGHC3\_v3 down by difficulty and ambiguity. Easy scenes remain high in RSR, but their SSR is lower, indicating that some residual errors are mask or same-instance failures after mostly correct target selection. Hard scenes remain below 35\% RSR, showing that blocker and visibility inference are still the main bottlenecks.

\begin{table}[t]
\centering
\caption{EGHC3\_v3 per-stratum performance over three splits. Each stratum contains 150 samples.}
\label{tab:strata}
\small
\begin{tabular}{lccc}
\toprule
Stratum & $N$ & RSR & SSR \\
\midrule
Easy, no ambiguity & 150 & 93.33 & 79.33 \\
Easy, ambiguity & 150 & 82.67 & 76.67 \\
Medium, no ambiguity & 150 & 60.00 & 56.67 \\
Medium, ambiguity & 150 & 63.33 & 62.67 \\
Hard, no ambiguity & 150 & 32.00 & 29.33 \\
Hard, ambiguity & 150 & 34.00 & 26.67 \\
\bottomrule
\end{tabular}
\end{table}

\subsection{Failure Analysis}

The stratum-level pattern suggests two different failure regimes. In easy scenes, EGHC usually resolves the requested target, so the remaining gap between RSR and SSR is consistent with instance-level mask ambiguity, missing masks, or masks that do not contain the selected Molmo point. In hard scenes, RSR itself is low; the dominant errors occur before mask evaluation, when the route or first-grasp ID does not match the physically relevant target/blocker relation.

Qualitative inspection should therefore be organized around two groups. Easy failures should show same-class or mask-IoU ambiguity. Hard failures should show wrong blocker selection, direct target selection when blocker removal was needed, or blocker selection when the target was directly accessible. This split keeps the discussion grounded: EGHC improves the decision contract, but it does not solve hard-scene visibility estimation.

\subsection{Statistical Verification}

The current tables report frozen point estimates. Before submission, we will add paired, scene-clustered bootstrap confidence intervals for EGHC3\_v3 versus S0 and for S1 versus EGHC3\_v3. The paired key will include split, subset, row index, scene ID, query object ID, and annotation text. Until those intervals are added, significance language is intentionally avoided.
